\documentclass[aspectratio=169,11pt]{beamer}
\usetheme{Madrid}
\usecolortheme{seahorse}

\usepackage{ctex}
\usepackage{listings}
\usepackage{xcolor}
\usepackage{graphicx}

\lstset{
    basicstyle=\ttfamily\footnotesize,
    keywordstyle=\color{blue},
    commentstyle=\color{gray},
    stringstyle=\color{orange},
    breaklines=true,
    frame=single,
    backgroundcolor=\color{gray!10}
}

\title{MultiTurn-RAG-CustomerService}
\subtitle{《中华人民共和国药典》智能问答系统架构梳理}
\author{Hermnd-Isabell}
\date{\today}

\begin{document}

% ========== 标题页 ==========
\begin{frame}
    \titlepage
\end{frame}

% ========== 目录 ==========
\begin{frame}{目录}
    \tableofcontents
\end{frame}

% ========== 1. 项目概述 ==========
\section{项目概述}

\begin{frame}{项目定位}
    \textbf{MultiTurn-RAG-CustomerService} 是一个基于《中华人民共和国药典》的\\中文智能问答系统。
    \vspace{0.5em}
    \begin{itemize}
        \item \textbf{R}etrieval: Elasticsearch + FAISS 双路检索
        \item \textbf{A}ugmented: 检索结果作为 LLM 上下文
        \item \textbf{G}eneration: OpenAI 兼容 API（默认 DeepSeek）生成回答
        \item \textbf{MultiTurn}: 支持多轮对话、指代消解、会话上下文
    \end{itemize}
\end{frame}

\begin{frame}{核心数据流}
    \centering
    \begin{tabular}{ll}
        \textbf{阶段} & \textbf{处理} \\
        \hline
        文档导入 & .docx $\rightarrow$ ES（整章）$\rightarrow$ FAISS（子段落） \\
        用户提问 & 查询改写 $\rightarrow$ 意图路由 $\rightarrow$ 检索决策 \\
        检索阶段 & ES 精确锁定（主路径）/ FAISS 向量检索（降级） \\
        生成阶段 & 上下文组装 $\rightarrow$ LLM 流式回答 \\
    \end{tabular}
\end{frame}

% ========== 2. 已实现功能 ==========
\section{已实现功能}

\begin{frame}{P0：基础架构}
    \begin{block}{双路存储}
        \begin{itemize}
            \item \textbf{Elasticsearch}：存储整章原文，doc\_id = 药品名
            \item \textbf{FAISS}：存储子段落向量（20857 条，384 维）
            \item 同一份文档存两份，目的不同
        \end{itemize}
    \end{block}
    \begin{block}{文档切分}
        \begin{itemize}
            \item 章节边界：检测 12pt 字体段落
            \item 子段落：正则 \texttt{(?:【|t)(.+?)(?:】)} 提取标题和内容
        \end{itemize}
    \end{block}
\end{frame}

\begin{frame}{P1：意图路由层}
    \begin{enumerate}
        \item \texttt{quick\_intent\_hint}：规则预筛
        \begin{itemize}
            \item obvious\_chitchat（闲聊短路）
            \item obvious\_pharmacy（含字段关键词）
            \item ambiguous（需 LLM 判断）
        \end{itemize}
        \item \texttt{classify\_pharmacy\_query}：LLM 二分类 good/bad
        \item \texttt{MedicineInfoStandardizer}：
        \begin{itemize}
            \item 从问句提取 \texttt{target\_fields}（如「性状」「鉴别」）
            \item 从问句提取 \texttt{target\_drug}（药品名）
        \end{itemize}
    \end{enumerate}
\end{frame}

\begin{frame}{P2：药品感知重排}
    全局 FAISS 检索 top-10 后，做双层打分重排：
    \vspace{0.5em}
    \begin{equation*}
        \text{total\_score} = \text{drug\_score} \times 2 + \text{field\_score} \times 1
    \end{equation*}
    \vspace{0.3em}
    \begin{itemize}
        \item \texttt{\_score\_result\_by\_drug\_name}：doc\_id 与目标药品名匹配（精确/包含）
        \item \texttt{\_score\_result\_by\_fields}：段落 title 与 target\_fields 匹配
        \item 匹配 $<$ 3 条时用向量相似度补足
    \end{itemize}
\end{frame}

\begin{frame}{P2.2：ES 精确锁定检索}
    当 \texttt{target\_drug} 非空且 ES 中存在时：
    \vspace{0.5em}
    \begin{itemize}
        \item \textbf{跳过}全库 FAISS 盲搜
        \item 直接 \texttt{es.get(index, id=target\_drug)} 取整章原文
        \item \texttt{extract\_subsections} 切分后按 \texttt{target\_fields} 过滤
        \item 匹配字段置顶，其余子段落补足 top\_k
    \end{itemize}
    \vspace{0.5em}
    \alert{效果}：「川射干的性状」直接命中【性状】段落，零漂移。
\end{frame}

\begin{frame}{P4.1：查询改写（Query Rewriting）}
    解决多轮对话中的\textbf{指代消解}问题：
    \vspace{0.5em}
    \begin{itemize}
        \item 规则预筛：已含药名 / 无指代词 $\rightarrow$ 跳过改写
        \item 指代词集合：\{他, 她, 它, 这个药, 该药, 刚才, 上面, 之前, ...\}
        \item 否则调用 LLM 结合最近 3 轮 history 做改写
    \end{itemize}
    \vspace{0.5em}
    \begin{exampleblock}{示例}
        Q1: 川射干的性状是什么？\\
        Q2: \textbf{他的副作用是什么？} $\xrightarrow{\text{改写}}$ \textbf{川射干的副作用是什么？}
    \end{exampleblock}
\end{frame}

\begin{frame}{P4.2：上下文感知检索}
    \begin{block}{\texttt{retrieve\_with\_context}}
        全库向量检索的\textbf{双保险}：
        \begin{enumerate}
            \item 基础检索 top-15（扩大候选池）
            \item 如果 \texttt{context\_drug} 非空：
            \begin{itemize}
                \item 用 \texttt{\_score\_result\_by\_drug\_name} 做药品名硬过滤
                \item 匹配结果置顶，不足时用基础结果补足
            \end{itemize}
            \item 任何异常降级为返回基础结果
        \end{enumerate}
    \end{block}
\end{frame}

\begin{frame}{P4.3：会话事实缓存}
    模块级 \texttt{\_session\_facts} 字典（Gradio 进程隔离 = 会话隔离）：
    \vspace{0.5em}
    \begin{itemize}
        \item \texttt{primary\_drug}：当前会话主要讨论药品
        \item \texttt{queried\_fields}：已查询字段集合（去重）
        \item \texttt{last\_intent}：上一轮意图
        \item \texttt{drug\_history}：会话中所有药品（按时间顺序）
    \end{itemize}
    \vspace{0.5em}
    作用：当前轮次提取不到药品名时，自动 fallback 到缓存值。
\end{frame}

\begin{frame}{测试覆盖}
    \begin{itemize}
        \item \texttt{tests/test\_config.py}：配置读取（17 个）
        \item \texttt{tests/test\_embed.py}：嵌入/检索/字段提取（44 个）
        \item \texttt{tests/test\_webrun.py}：slow\_echo 分支/重排/改写/缓存（33 个）
        \item \texttt{tests/test\_drug\_aware\_rerank.py}：药品名打分/提取（18 个）
    \end{itemize}
    \vspace{0.5em}
    \centering
    \textbf{总计：128 个测试全部通过}
\end{frame}

% ========== 3. 实现机制详解 ==========
\section{实现机制详解}

\begin{frame}{检索决策流程图}
    \centering
    \small
    \begin{tabular}{ll}
        \textbf{条件} & \textbf{路径} \\
        \hline
        target\_drug $\neq$ None \&\& drug\_exists & \textbf{主路径}：ES 精确锁定 \\
        target\_drug = None \&\& session\_drug $\neq$ None & \textbf{降级路径}：retrieve\_with\_context \\
        target\_drug = None \&\& session\_drug = None & \textbf{降级路径}：FAISS 全库 top-3 \\
        全库未命中目标药品 & \textbf{兜底}：retrieve\_vector\_and\_text\_for\_drug \\
    \end{tabular}
\end{frame}

\begin{frame}{关键模块关系}
    \centering
    \begin{tabular}{lll}
        \textbf{模块} & \textbf{职责} & \textbf{关键函数} \\
        \hline
        config.py & 全局配置单例 & \texttt{config.ES\_INDEX} \\
        embed.py & 向量/ES/LLM 工具 & \texttt{retrieve\_with\_context} \\
        webrun.py & Gradio UI + 主流程 & \texttt{slow\_echo} \\
        tests/ & 回归测试 & pytest + MagicMock \\
    \end{tabular}
\end{frame}

% ========== 4. 当前问题 ==========
\section{当前问题与改进方向}

\begin{frame}{已知问题（1/2）}
    \begin{alertblock}{结构性问题}
        \begin{itemize}
            \item \textbf{向量语义漂移}：FAISS 只编码子段落\textit{内容}，不含药品名\\
            $\Rightarrow$ 「川射干的性状」编码后失去「川射干」信号
            \item \textbf{切分规则硬编码}：依赖 12pt 字体 + 固定正则\\
            $\Rightarrow$ 非标准 .docx 无法正确分章
            \item \textbf{模型不可切换}：Sentence-Transformer 固定加载\\
            $\Rightarrow$ 无法根据场景换用更强模型
        \end{itemize}
    \end{alertblock}
\end{frame}

\begin{frame}{已知问题（2/2）}
    \begin{alertblock}{工程性问题}
        \begin{itemize}
            \item \textbf{LLM 依赖}：改写/分类/字段提取均走外部 API\\
            $\Rightarrow$ 延迟高、有失败风险、成本高
            \item \textbf{会话状态易失}：\texttt{\_session\_facts} 是模块级变量\\
            $\Rightarrow$ Gradio 进程重启即丢失
            \item \textbf{无对话持久化}：没有数据库记录历史问答
            \item \textbf{Windows-only}：启动脚本为 .bat，跨平台性差
            \item \textbf{无权限管理}：配置页密码明文，无用户隔离
        \end{itemize}
    \end{alertblock}
\end{frame}

\begin{frame}{改进方向}
    \begin{exampleblock}{短期}
        \begin{itemize}
            \item 向量库编码时 prepend 药品名（\texttt{[川射干] 本品为...}）
            \item 引入本地小模型做意图分类（减少 LLM 调用）
            \item 添加 SQLite 对话历史持久化
        \end{itemize}
    \end{exampleblock}
    \begin{exampleblock}{中长期}
        \begin{itemize}
            \item 支持多药品对比查询
            \item 动态 field\_list 扩展（用户自定义字段）
            \item 跨平台启动脚本（sh + Docker）
            \item 引入重排序模型（Cross-Encoder）替代规则打分
        \end{itemize}
    \end{exampleblock}
\end{frame}

% ========== 结束页 ==========
\begin{frame}
    \centering
    \Huge 谢谢！
    \vspace{1em}
    \\
    \normalsize
    GitHub: \url{https://github.com/Hermnd-Isabell/MultiTurn-RAG-CustomerService}
\end{frame}

\end{document}
